\begin{abstract}
Long-term memory lets an agent answer questions about past conversations. For text memory, recent studies find that simply retrieving the raw history by similarity matches or beats dedicated memory systems, and we confirm that it comes within one point of reading the whole history. We show that this result does not carry over to multimodal memory. When the evidence lies in images, many questions ask about a set of past items, such as every receipt from a trip, and ranking items one by one rarely returns the whole set, so retrieval falls well short of reading the evidence. We present Visibility-Gated Set Reading (\vgs{}), a training-free memory reader that separates navigation from perception: the memory keeps raw rounds and writes one short index line per round, and a model reads the index and returns the evidence set, choosing its size. Because a model can only find what the index shows, \vgs{} estimates this index visibility and widens its reading when the index hides the evidence. On three multimodal benchmarks, \vgs{} outperforms top-40 retrieval and six recent memory systems with a lightweight reader on all three benchmarks and with a strong reader wherever the index exposes the evidence. Ablations show that reading the index rather than ranking by similarity recovers the complete sets, and that the gate picks the better reading on every benchmark.
\end{abstract}
